% !TEX program = xelatex
% ============================================================
%  爱因斯坦-德哈斯效应相关文献 中英对照集  —  公共导言
% ============================================================
\usepackage{xeCJK}
\usepackage{fontspec}
\usepackage{geometry}
\usepackage{setspace}
\usepackage{enumitem}
\usepackage{amsmath,amssymb}
\usepackage{graphicx}
\usepackage{adjustbox}
\usepackage{booktabs}
\usepackage{longtable}
\usepackage{array}
\usepackage{fancyhdr}
\usepackage{titlesec}
\usepackage{titletoc}
\usepackage[most]{tcolorbox}
\usepackage{needspace}
\usepackage{footmisc}
\usepackage{xcolor}
\usepackage{hyperref}
\AtBeginDocument{\renewcommand{\contentsname}{目录}}

\geometry{a4paper,left=2.4cm,right=2.4cm,top=2.6cm,bottom=2.6cm}

% ---------- 防止内容溢出/重叠（务必保留） ----------
% 让 LaTeX 放宽断行，避免行溢出伸进页边距与相邻内容重叠
\setlength{\emergencystretch}{3em}
\tolerance=2000
\hbadness=3000
% 禁止孤行寡行（段落首尾行悬空，配合图块防重叠）
\clubpenalty=10000
\widowpenalty=10000
\displaywidowpenalty=10000

% ---------- 字体 ----------
\setmainfont{Times New Roman}
\setsansfont{Arial}
\setCJKmainfont{SimSun}[BoldFont=SimHei,ItalicFont=FandolKai]
\setCJKsansfont{SimHei}
\setCJKfamilyfont{kai}{FandolKai}

\newcommand{\kai}{\CJKfamily{kai}}

% ---------- 颜色 ----------
\definecolor{enColor}{RGB}{20,20,20}
\definecolor{zhColor}{RGB}{10,40,90}
\definecolor{termColor}{RGB}{150,20,20}
\definecolor{capColor}{RGB}{90,90,90}
\definecolor{ruleColor}{RGB}{180,180,180}

% ---------- 标题样式 ----------
\renewcommand{\chaptername}{}
\renewcommand{\figurename}{图}
\renewcommand{\tablename}{表}
\titleformat{\chapter}[display]
  {\bfseries\sffamily\Large}{}{0pt}{\Large}
% 顶距不得为负：负值会把章名顶进页眉造成重叠
\titlespacing*{\chapter}{0pt}{4pt}{18pt}
% 目录只显示到章（每篇文献一条），不列小节
\setcounter{tocdepth}{0}
% 目录中章条目带编号
\titlecontents{chapter}[0em]
  {\addvspace{8pt}\bfseries\sffamily}
  {\thecontentslabel\quad}{}
  {\hfill\contentspage}

\titleformat{\section}{\bfseries\sffamily\large}{\thesection}{0.8em}{}
\titleformat{\subsection}{\bfseries\sffamily\normalsize}{\thesubsection}{0.8em}{}

% ---------- 页眉页脚 ----------
\pagestyle{fancy}
\fancyhf{}
\fancyhead[L]{}
\fancyhead[R]{\small\sffamily\leftmark}
\fancyfoot[C]{\small\thepage}
\renewcommand{\headrulewidth}{0.3pt}
\renewcommand{\footrulewidth}{0pt}
\renewcommand{\chaptermark}[1]{\markboth{#1}{}}

% ---------- 段落环境 ----------
\newenvironment{enpara}%
  {\par\smallskip\begingroup\color{enColor}\setstretch{1.05}\noindent}%
  {\par\endgroup\smallskip}
\newenvironment{zhpara}%
  {\par\begingroup\color{zhColor}\setstretch{1.35}\noindent}%
  {\par\endgroup\medskip}

% 术语（首次出现加脚注），正文中把术语标红
\newcommand{\term}[2]{\textcolor{termColor}{\textbf{#1}}%
  \footnote{\textcolor{termColor}{\textbf{#1}}：\kai #2}}

% 图表标题
\newcommand{\figcap}[2]{\par{\small\color{capColor}\noindent #1}\par
  {\footnotesize\color{capColor}\noindent\itshape #2}\par\smallskip}

% 分隔线
\newcommand{\pairrule}{\par\vspace{0.1em}{\color{ruleColor}\hrule height 0.4pt}\vspace{0.2em}}

% 摘要标题（中英标注）
\newcommand{\abstracthead}{\par\medskip\noindent
  {\sffamily\bfseries\large\color{zhColor}摘要\ \textnormal{\itshape /\ Abstract}}\par\smallskip
  {\color{ruleColor}\hrule height 0.6pt}\vspace{0.3em}}
